\subsubsection{Filtering and Joins}

The foundation. Every later pattern assumes you can already restrict rows and
combine tables.

\paragraph{WHERE vs HAVING.} \texttt{WHERE} filters individual rows before
grouping. \texttt{HAVING} filters whole groups after aggregation and is the
only place an aggregate predicate is legal. A filter on a plain
(non-aggregated) column belongs in \texttt{WHERE} so rows are dropped before
the grouping work happens. Reserve \texttt{HAVING} for conditions like
\texttt{COUNT(*) > 5}.

\begin{lstlisting}[style=sql, caption={WHERE prunes rows, HAVING prunes groups}]
SELECT department, COUNT(*) AS headcount
FROM employees
WHERE active = true          -- row filter, before grouping
GROUP BY department
HAVING COUNT(*) > 5;         -- group filter, on an aggregate
\end{lstlisting}

\vspace{-1ex}

\newpage

\paragraph{Joins.} The recognition cue is which side you want to keep when a
match is missing.
\begin{itemize}[label=$\circ$]
  \item \textbf{INNER:} rows with a match on both sides.
  \item \textbf{LEFT / RIGHT:} keep all rows from one side, \texttt{NULL}-fill the other.
  \item \textbf{FULL OUTER:} keep unmatched rows from both sides.
  \item \textbf{CROSS:} Cartesian product, every left row against every right row.
  \item \textbf{SELF:} a table joined to itself under an alias (employee to manager, row to prior row).
  \item \textbf{Semi-join} (\texttt{EXISTS}, \texttt{IN}): keep left rows with at least one match. No row duplication, no right-side columns.
  \item \textbf{Anti-join} (\texttt{NOT EXISTS}, or \texttt{LEFT JOIN ... WHERE right.key IS NULL}): keep left rows with no match.
\end{itemize}

Semi- and anti-joins filter the left table on whether a match exists, without
pulling right-side columns or duplicating rows. Prefer \texttt{NOT EXISTS} over
\texttt{NOT IN} for the anti-join, since a \texttt{NULL} on the right poisons
\texttt{NOT IN} to zero rows.

\begin{lstlisting}[style=sql, caption={Semi-join keeps matches, anti-join keeps non-matches}]
-- Semi-join: customers with at least one order
SELECT c.customer_id, c.name
FROM customers c
WHERE EXISTS (SELECT 1 FROM orders o
              WHERE o.customer_id = c.customer_id);

-- Anti-join: customers with no orders
SELECT c.customer_id, c.name
FROM customers c
WHERE NOT EXISTS (SELECT 1 FROM orders o
                  WHERE o.customer_id = c.customer_id);

-- Anti-join, LEFT JOIN form: unmatched rows are the NULL-filled ones
SELECT c.customer_id, c.name
FROM customers c
LEFT JOIN orders o ON o.customer_id = c.customer_id
WHERE o.customer_id IS NULL;
\end{lstlisting}